\documentclass{article}
\usepackage{amsmath, amssymb}

\begin{document}

\section*{Pushing--Medium Gravity v2: Complete Unified Formulation with Numerical Constraints}

This document presents the full gravitational sector of the Pushing--Medium (PM)
model, including static fields, dynamic equations, the Lagrangian formulation,
ray and particle kinematics, the effective metric derived from Fermat’s
principle, and constraints from numerical tests. Version 2 incorporates
convergence, iterative-deflection, and weak-bending constraints.

The PM gravitational field is encoded in two fields:



\[
\phi(\mathbf r,t) \equiv \ln n(\mathbf r,t),
\qquad
\mathbf u(\mathbf r,t) \equiv \text{gravitational flow field}.
\]



The scalar field $\phi$ determines the refractive index $n$, while the flow
field $\mathbf u$ encodes moving-mass effects.

\section*{1. Unified Source Structure}

\subsection*{1.1 Scalar Sources}



\[
S_\phi = 4\pi G_{\mathrm{eff}}\,\rho_M(\mathbf r,t),
\]



where $\rho_M$ is mass density.

\subsection*{1.2 Vector Sources}



\[
\mathbf S_u = A_M\,\mathbf J_M(\mathbf r,t),
\]



where $\mathbf J_M$ is mass current density.

\section*{2. Static PM Gravity}

\subsection*{2.1 Scalar Index Field}

In the static limit ($\partial_t \phi = 0$, $\mathbf u = 0$):



\[
\nabla^2 \phi(\mathbf r)
=
-4\pi G_{\mathrm{eff}}\,\rho_M(\mathbf r).
\]



The refractive index is:



\[
n(\mathbf r) = \exp(\phi(\mathbf r)).
\]



\subsection*{2.2 Point Mass}

For a point mass $M$:



\[
\phi(r) = \frac{\mu_G M}{r},
\qquad
n(r) = 1 + \frac{\mu_G M}{r} + \mathcal{O}(r^{-2}).
\]



\subsection*{2.3 Continuous Mass Distribution}



\[
\phi(\mathbf r)
=
\mu_G
\int
\frac{\rho_M(\mathbf r')}{|\mathbf r - \mathbf r'|}
\,d^3r'.
\]



\subsection*{2.4 Static Flow Field}



\[
\nabla^2 \mathbf u(\mathbf r)
=
- A_M\,\mathbf J_M(\mathbf r).
\]



For a moving point mass:



\[
\mathbf u(\mathbf r)
=
A_M M
\frac{\mathbf v}{|\mathbf r - \mathbf r_M|}.
\]



\section*{3. Dynamic PM-Fluid Equations}

The PM medium behaves as a compressible, advective fluid with finite
propagation speeds $c_\phi$ and $c_u$.

\subsection*{3.1 Scalar Field Equation}



\[
\frac{\partial \phi}{\partial t}
+ \nabla\cdot(\phi\,\mathbf u)
=
S_\phi
+
\frac{1}{c_\phi^2}\frac{\partial^2 \phi}{\partial t^2}.
\]



\subsection*{3.2 Vector Flow Equation}



\[
\frac{\partial \mathbf u}{\partial t}
+ (\mathbf u\cdot\nabla)\mathbf u
=
-\nabla \phi
+ \mathbf S_u
+ \frac{1}{c_u^2}\frac{\partial^2 \mathbf u}{\partial t^2}.
\]



These reduce to the static PM gravity equations when $\partial_t = 0$ and
$\mathbf u = 0$.

\section*{4. Lagrangian Formulation}

The PM-fluid equations arise from the Lagrangian density:



\[
\mathcal{L}
=
\frac{1}{2c_\phi^2}(\partial_t \phi)^2
- \frac{1}{2}|\nabla \phi|^2
+
\frac{1}{2c_u^2}|\partial_t \mathbf u|^2
- \frac{1}{2}|\nabla \mathbf u|^2
+
\phi\,\nabla\cdot(\phi\mathbf u)
+
\phi\,S_\phi
+
\mathbf u\cdot\mathbf S_u.
\]



The terms represent:

\begin{itemize}
    \item scalar and vector kinetic energies,
    \item spatial gradient energies,
    \item advection/compressibility,
    \item unified source couplings.
\end{itemize}

\section*{5. Ray and Particle Kinematics}

Rays propagate according to:



\[
\frac{d\mathbf r}{ds}
= c\,\hat{\mathbf k}
+ k_{\mathrm{Fizeau}}\,\mathbf u(\mathbf r,t),
\]





\[
\frac{d\hat{\mathbf k}}{ds}
= \nabla_\perp \phi(\mathbf r,t)
+ \mathcal{F}\big(\hat{\mathbf k}, \mathbf u, \nabla \mathbf u\big).
\]



Empirically, $k_{\mathrm{Fizeau}} = 1$.

\section*{6. Effective Metric from Fermat’s Principle}

For a moving medium with index $n = e^\phi$ and flow $\mathbf u$, the optical
metric is:



\[
ds_{\mathrm{opt}}^2
=
-\frac{c^2}{n^2}\,dt^2
+
\left(d\mathbf r - \mathbf u\,dt\right)^2.
\]



Expanding to first order in $\phi$ and $\mathbf u$:



\[
ds^2
=
-(1 + 2\phi)c^2 dt^2
- 2\,\mathbf u\cdot d\mathbf r\,dt
+
(1 - 2\phi)\,d\mathbf r^2.
\]



\subsection*{6.1 PPN Parameters}



\[
g_{tt} = -(1 + 2\phi),
\qquad
g_{rr} = (1 - 2\phi),
\]



yielding:



\[
\gamma = 1,
\qquad
\beta = 1.
\]



\section*{7. Numerical Constraints (v2 Additions)}

\subsection*{7.1 Weak/Semistrong Deflection}

Tests show:



\[
\alpha(b) \propto \frac{1}{b},
\]



with log–log slope $k \approx 1$ and $\alpha b$ approximately constant.

This confirms:

\begin{itemize}
    \item linearity of the scalar field in this regime,
    \item no saturation or super-divergence,
    \item consistency with the PM-fluid v1 equations.
\end{itemize}

\subsection*{7.2 Deflection Convergence}

Numerical convergence tests show:

\begin{itemize}
    \item stable refinement under decreasing step size,
    \item no stiffness or divergence in the ray equation,
    \item no need for viscosity-like or higher-derivative stabilizers.
\end{itemize}

Thus the PM-fluid v1 equations require no additional regularization.

\subsection*{7.3 Iterative Deflection}

Iterative-deflection tests show:

\begin{itemize}
    \item fixed-point iteration converges without damping,
    \item no evidence of nonlinear feedback,
    \item no need for nonlinear advection or potential terms.
\end{itemize}

\subsection*{7.4 Weak-Bending Numeric}

Weak-bending tests confirm:



\[
\alpha_{\mathrm{PM}}(b)
=
C\,\frac{1}{b}
\quad\text{with constant } C,
\]



matching the PM index field and the effective metric.

This validates:

\begin{itemize}
    \item the linear scalar field equation,
    \item the PM ray equation,
    \item the weak-field metric structure.
\end{itemize}

\section*{8. Critical Compression and Matter Stability}

In the high-compression regime the PM scalar field $\phi$ develops
non-linear self-interactions that set an upper stability limit for matter
configurations.  The deformation energy density (see
\texttt{hamiltonian-formulation-v1.tex}, \S7) is:


\[
U(\phi) = \varepsilon_0\!\left(\phi^2 - \frac{\phi^3}{3}\right),
\qquad
U''(\phi) = 2\varepsilon_0(1-\phi).
\]


The critical compression $\phi_{\mathrm{crit}} = 1$ is where $U'' = 0$.  Using
the PM density mapping $\rho = \rho_{\mathrm{nuc}}e^\phi$ and ultra-relativistic
EOS $P = \rho c^2/3$:


\[
\phi_{\mathrm{crit}} = 1,
\qquad
n_{\mathrm{crit}} = e,
\qquad
\rho_{\mathrm{crit}} \approx 6.25\times10^{17}\;\mathrm{kg\,m^{-3}},
\qquad
P_{\mathrm{crit}} \approx 1.87\times10^{34}\;\mathrm{Pa}.
\]


This places the PM matter/energy transition at $\approx 2.7\,\rho_{\mathrm{nuc}}$,
within the expected nuclear-failure band and consistent with the broad
Einstein--Oppenheimer--Volkoff picture.  Full documentation, implementation,
and tests are in \texttt{critical\_state\_computation.md},
\texttt{src/pushing\_medium/critical\_state.py}, and
\texttt{tests/test\_critical\_state.py}.

\section*{9. Summary}

The PM gravitational sector consists of:

\begin{itemize}
    \item a scalar index field sourced by mass density,
    \item a flow field sourced by mass currents,
    \item dynamic PM-fluid equations with advection and finite propagation,
    \item a Lagrangian generating the full dynamics,
    \item an effective metric matching weak-field GR,
    \item numerical validation of linearity, convergence, and weak/semistrong bending,
    \item a deformation-energy stability criterion fixing the matter/energy transition at $\phi_{\mathrm{crit}} = 1$.
\end{itemize}

This v2 formulation incorporates all current numerical constraints and provides
a complete, self-consistent gravitational theory within the Pushing--Medium
framework.

\end{document}

